% Luc Destombe --- licence CC BY-NC-SA 4.0
% https://creativecommons.org/licenses/by-nc-sa/4.0/deed.fr
% Fichier autonome : compiler avec pdflatex, deux fois (signets, nombre de pages).
\documentclass[11pt,a4paper]{article}
\makeatletter
\newif\ifeleve
\elevefalse

\newif\iflycee@palatino
\lycee@palatinotrue

\RequirePackage[utf8]{inputenc}
\RequirePackage[T1]{fontenc}
\RequirePackage[french]{babel}
\RequirePackage{etoolbox}

\newcommand{\ShorthandsOff}{\@ifundefined{shorthandoff}{}{\shorthandoff{;:!?}}}
\newcommand{\ShorthandsOn}{\@ifundefined{shorthandon}{}{\shorthandon{;:!?}}}
\ShorthandsOff
\AtBeginDocument{\ShorthandsOn}
\AtBeginEnvironment{tikzpicture}{\ShorthandsOff}

\RequirePackage{geometry}
\RequirePackage{fancyhdr}
\RequirePackage{lastpage}
\RequirePackage{multicol}
\RequirePackage{array}
\RequirePackage{enumitem}
\RequirePackage{xcolor}
\RequirePackage{needspace}

\RequirePackage{amsmath,amssymb}
\RequirePackage{mathpazo}
\DeclareMathSizes{10.95}{10.95}{8}{5.5}
\begingroup\catcode`\_=\active \catcode`\^=\active
\gdef_#1{\sb{\scriptscriptstyle #1}}
\gdef^#1{\sp{\scriptscriptstyle\lycee@moinsepais #1}}
\endgroup
\newcommand\lycee@moins{\mathbin{%
  \setbox\z@\hbox{$\m@th\scriptscriptstyle\lycee@moinsorig$}%
  \dimen@=.55\wd\z@ \advance\dimen@-.5pt
  \hbox to.75\wd\z@{\hss$\m@th\scriptscriptstyle\vcenter{\hbox{%
    \tikz\draw[line width=.5pt,line cap=round](0,0)--(\the\dimen@,0);}}$\hss}}}
\begingroup\lccode`\~=`\- \lowercase{\endgroup\let~}\lycee@moins
\newcommand\lycee@moinsepais{\mathcode`\-="8000 }
\AtBeginDocument{\mathchardef\lycee@moinsorig=\mathcode`\-\relax}
\AtBeginDocument{\catcode`\_=12 \mathcode`\_="8000
  \catcode`\^=12 \mathcode`\^="8000 }
\newcommand\lycee@exposants{%
  \fontdimen13\textfont2=.48\fontdimen6\textfont2
  \fontdimen14\textfont2=.48\fontdimen6\textfont2
  \fontdimen15\textfont2=.48\fontdimen6\textfont2
  \fontdimen13\scriptfont2=.48\fontdimen6\scriptfont2
  \fontdimen14\scriptfont2=.48\fontdimen6\scriptfont2
  \fontdimen15\scriptfont2=.48\fontdimen6\scriptfont2}
\every@math@size\expandafter{\the\every@math@size\lycee@exposants}
\newlength{\retraitcalculs}
\setlength{\retraitcalculs}{2em}

\RequirePackage{tikz}
\usetikzlibrary{arrows.meta,calc,babel,shapes.misc}
\RequirePackage[most]{tcolorbox}
\tcbuselibrary{skins,breakable}

\definecolor{coulDef}{HTML}{1F4E79}
\definecolor{coulProp}{HTML}{7030A0}
\definecolor{coulRem}{HTML}{C55A11}
\definecolor{coulEx}{HTML}{2E7D32}
\definecolor{coulExo}{HTML}{B03A2E}
\definecolor{coulPart}{HTML}{1F4E79}
\definecolor{coulMeth}{HTML}{1B6E4F}
\definecolor{coulCourbe}{HTML}{0B5ED7}
\definecolor{coulCourbeB}{HTML}{C00000}
\definecolor{coulCourbeC}{HTML}{2E7D32}
\definecolor{coulGrille}{HTML}{8C8C8C}
\definecolor{coulRep}{HTML}{1B6E4F}
\definecolor{coulBar}{HTML}{2E7D32}

\newcommand{\R}{\mathbb{R}}
\newcommand{\Rs}{\mathbb{R}^{*}}
\renewcommand{\le}{\leqslant}
\renewcommand{\ge}{\geqslant}
\renewcommand{\approx}{\simeq}
\newcommand{\vect}[1]{\overrightarrow{#1}}
\newcommand{\norme}[1]{\left\lVert #1 \right\rVert}
\newcommand{\intervalle}[2]{\left[#1\,;#2\right]}
\RequirePackage{cancel}
\renewcommand{\CancelColor}{\color{coulExo}}
\newcommand{\fois}[1]{\times{\color{coulExo}#1}}
\newcommand{\barre}[1]{\cancel{#1}}

\tikzset{grille/.style={coulGrille,line width=0.4pt}}
\newcommand{\quadrillage}[4]{\draw[grille,step=1] (#1,#2) grid (#3,#4);}
\tikzset{
  croix/.style={cross out,draw,line width=0.6pt,inner sep=0pt,minimum size=4pt},
  croixdroite/.style={croix,rotate=45,line width=0.8pt},
}
\newcommand{\pt}[3]{\path (#1) node[croix]{} node[#2,font=\small,inner sep=2.5pt]{$#3$};}

\newcommand{\lycee@repere}[4]{%
  \draw[grille,step=1] (#1,#2) grid (#3,#4);
  \draw[-{Stealth[length=2mm]},thick] (#1,0)--({#3+0.4},0) node[below left=-1pt]{$x$};
  \draw[-{Stealth[length=2mm]},thick] (0,#2)--(0,{#4+0.4}) node[below left=-1pt]{$y$};
  \node[below left,font=\scriptsize,inner sep=1.5pt] at (0,0) {$0$};}
\newcommand{\lycee@gradx}[1]{\foreach \k/\l in {#1}{%
  \draw (\k,0.07)--(\k,-0.07) node[below,font=\scriptsize,inner sep=1.5pt]{$\l$};}}
\newcommand{\lycee@grady}[1]{\foreach \k/\l in {#1}{%
  \draw (0.07,\k)--(-0.07,\k) node[left,font=\scriptsize,inner sep=1.5pt]{$\l$};}}
\newcommand{\lycee@intff}[2]{\left[#1\,;#2\right]}
\newcommand{\lycee@intfo}[2]{\left[#1\,;#2\right[}
\newcommand{\lycee@intof}[2]{\left]#1\,;#2\right]}
\newcommand{\lycee@intoo}[2]{\left]#1\,;#2\right[}
\newcommand{\lycee@ens}[1]{\left\{#1\right\}}
\AtBeginDocument{%
  \providecommand{\repere}{\lycee@repere}%
  \providecommand{\gradx}{\lycee@gradx}%
  \providecommand{\grady}{\lycee@grady}%
  \providecommand{\Cf}{\mathcal{C}\sb{\scriptscriptstyle f}}%
  \providecommand{\Cg}{\mathcal{C}\sb{\scriptscriptstyle g}}%
  \providecommand{\intff}{\lycee@intff}%
  \providecommand{\intfo}{\lycee@intfo}%
  \providecommand{\intof}{\lycee@intof}%
  \providecommand{\intoo}{\lycee@intoo}%
  \providecommand{\ens}{\lycee@ens}}

\newlist{questions}{enumerate}{3}
\setlist[questions,1]{label=\arabic*\textdegree),leftmargin=*,itemsep=2pt,topsep=3pt}
\setlist[questions,2]{label=\alph*),leftmargin=*,itemsep=1pt,topsep=2pt}
\setlist[questions,3]{label=\roman*),leftmargin=*,itemsep=1pt,topsep=1pt}
\setlist[itemize]{label=\textbullet,leftmargin=*,itemsep=2pt,topsep=3pt}

\newcommand{\besoinplace}[1]{\par\penalty\@M\begingroup
  \dimen@=#1\relax\dimen@ii=\pagegoal\advance\dimen@ii-\pagetotal
  \ifdim\dimen@>\dimen@ii\ifdim\dimen@ii>\z@\vfil\fi\break\fi\endgroup}

\newcommand{\lycee@pied}{\thepage/\pageref{LastPage}}
\newcommand{\entete}[2]{%
  \pagestyle{fancy}\fancyhf{}%
  \fancyhead[L]{\footnotesize\color{coulPart}\textbf{#1}}%
  \fancyhead[R]{\footnotesize\color{coulPart}\textbf{#2}}%
  \fancyfoot[C]{\footnotesize\lycee@pied}%
  \renewcommand{\headrulewidth}{0.4pt}%
  \renewcommand{\headrule}{\hbox to\headwidth{%
    \color{coulPart!40}\leaders\hrule height \headrulewidth\hfill}}}

\RequirePackage{ccicons}
\newcommand{\lycee@licence}{{\scriptsize\color{gray}%
  \copyright\ Luc Destombe --- \ccbyncsa\ CC BY-NC-SA 4.0}}
\newif\iflycee@licence \lycee@licencetrue
\newcommand{\sanslicence}{\lycee@licencefalse}
\AtBeginDocument{\iflycee@licence\AddToHookNext{shipout/foreground}{%
  \put(0,0){\rlap{\kern\dimexpr1in+\hoffset+\oddsidemargin+\textwidth\relax
    \llap{\raisebox{-\dimexpr1in+\voffset+\topmargin+\headheight+\headsep
      +\textheight+\footskip\relax}{\lycee@licence}}}}}\fi}

\newcommand{\formulesserrees}{%
  \setlength{\abovedisplayskip}{3pt}\setlength{\belowdisplayskip}{3pt}%
  \setlength{\abovedisplayshortskip}{2pt}\setlength{\belowdisplayshortskip}{3pt}}

\setlength{\parindent}{0pt}

\RequirePackage[implicit=false,hidelinks,pdfencoding=auto,psdextra,bookmarksopen,
  bookmarksopenlevel=1,pdfcreator={},pdfproducer={}]{hyperref}
\RequirePackage{bookmark}
\hypersetup{pdfauthor={Luc Destombe},pdfkeywords={CC BY-NC-SA 4.0}}
\newcounter{lycee@signet}
\newcommand{\signet}[3][]{\stepcounter{lycee@signet}%
  \edef\lycee@dest{\if\relax\detokenize{#1}\relax
    signet.\arabic{lycee@signet}\else#1\fi}%
  \hypertarget{\lycee@dest}{}\bookmark[level=#2,dest=\lycee@dest]{#3}}
\pdfstringdefDisableCommands{%
  \def\R{R}\def\vect#1{#1}\def\Cf{Cf}\def\Cg{Cg}\def\fois#1{×#1}%
  \def\barre#1{#1}\def\intervalle#1#2{[#1;#2]}\def\intff#1#2{[#1;#2]}%
  \def\textsuperscript#1{#1}\def\dfrac#1#2{#1/#2}\def\frac#1#2{#1/#2}%
  \def\vec#1{#1⃗}}%
\geometry{top=2cm,bottom=1cm,left=1cm,right=1cm,headheight=14pt,footskip=0.6cm}

\RequirePackage{pgfplots}
\pgfplotsset{compat=1.18}
\RequirePackage{tkz-tab}

\newcounter{partiecours}
\renewcommand{\thepartiecours}{\Roman{partiecours}}
\newcommand{\partiecours}[1]{%
  \refstepcounter{partiecours}%
  \Needspace*{8\baselineskip}%
  \bigskip
  \noindent\begin{tcolorbox}[enhanced,colback=coulPart,colframe=coulPart,
      boxrule=0pt,arc=2pt,left=8pt,right=8pt,top=4pt,bottom=4pt,
      before skip=6pt,after skip=10pt]
    \color{white}\bfseries\large \thepartiecours{} --- #1%
    \signet[partie-\arabic{partiecours}]{1}{\thepartiecours{} --- #1}
  \end{tcolorbox}%
  \addcontentsline{toc}{section}{\thepartiecours{} --- #1}%
}

\newtcolorbox{boitecours}[3][]{%
  enhanced,breakable,
  colback=white,colframe=#2,coltitle=white,
  fonttitle=\bfseries,
  attach boxed title to top left={xshift=8pt,yshift=-\tcboxedtitleheight/2},
  boxed title style={colback=#2,arc=2pt,boxrule=0pt},
  boxrule=0.9pt,arc=2pt,
  left=8pt,right=8pt,top=8pt,bottom=6pt,
  title={#3},#1}

\newcounter{methode}
\newenvironment{definitionbox}[1][Définition]%
  {\begin{boitecours}[phantom={\signet{2}{#1}}]{coulDef}{#1}}{\end{boitecours}}
\newenvironment{proprietebox}[1][Propriété]%
  {\begin{boitecours}[phantom={\signet{2}{#1}}]{coulProp}{#1}}{\end{boitecours}}
\newenvironment{methodebox}[1][Méthode]{\stepcounter{methode}%
  \begin{boitecours}[phantom={\signet[methode-\arabic{methode}]{2}{#1}}]{coulMeth}{#1}}%
  {\end{boitecours}}

\newcommand{\etape}[1]{\par\smallskip\textbf{\color{coulMeth}#1}\par\nobreak\smallskip}
\newcommand{\enonce}[1]{\emph{#1}\par\smallskip}
\newcommand{\reponse}[1]{\par\smallskip
  {\footnotesize\itshape\color{black!55}Réponses : #1}}
\newcommand{\demo}[1]{\textbf{\color{coulMeth}Démonstration.} #1}
\newcommand{\afaire}[1]{%
  \par\smallskip
  \noindent{\small\color{coulExo}$\blacktriangleright$\ \textbf{À faire :}
    fiche d'exercices du chapitre, #1.}\par\smallskip}

\newenvironment{calculs}{\setlength{\jot}{5pt}\@fleqntrue\@mathmargin\retraitcalculs
  \setlength{\abovedisplayskip}{4pt}\setlength{\belowdisplayskip}{4pt}%
  \setlength{\abovedisplayshortskip}{4pt}\setlength{\belowdisplayshortskip}{4pt}%
  \csname alignat*\endcsname{2}}%
  {\csname endalignat*\endcsname}
\newcommand{\rem}[1]{\text{\small\itshape\color{black!60}#1}}
\newcommand{\explique}[2]{\par\smallskip\noindent
  \begin{minipage}[t]{0.57\linewidth}\raggedright #1\end{minipage}\hfill
  \begin{minipage}[t]{0.40\linewidth}\raggedright\leavevmode
    {\small\itshape\color{black!60}#2}\end{minipage}\par}
\newcommand{\cas}[1]{\par\smallskip\noindent\textbf{\color{coulExo}#1}\nobreak}
\newcommand{\calc}[1]{\par\smallskip\textbf{\color{coulExo}#1}\par\nobreak}

\newtcolorbox{remarquebox}[1][Remarque]{%
  enhanced,breakable,colback=white,colframe=coulRem,
  boxrule=0pt,leftrule=2.5pt,arc=0pt,outer arc=0pt,
  left=8pt,right=6pt,top=5pt,bottom=5pt,
  before upper={\textbf{\color{coulRem}#1}\par\smallskip}}

\newtcolorbox{exemplebox}[1][Exemple]{%
  enhanced,breakable,colback=white,colframe=coulEx,
  boxrule=0pt,leftrule=2.5pt,arc=0pt,outer arc=0pt,
  left=8pt,right=6pt,top=5pt,bottom=5pt,
  before upper={\textbf{\color{coulEx}#1}\par\smallskip}}

\pgfplotsset{
  repere/.style={
    axis lines=middle,
    axis line style={-{Stealth[length=2.4mm]},thick},
    every axis x label/.style={at={(ticklabel* cs:1.0)},anchor=west},
    every axis y label/.style={at={(ticklabel* cs:1.0)},anchor=south},
    label style={font=\footnotesize},
    tick label style={font=\scriptsize},
    grid=both,
    major grid style={grille},
    minor tick num=0,
    tick align=inside,
    clip mode=individual,
  },
}
\tikzset{
  courbe/.style={coulCourbe,line width=1.1pt,smooth},
  courbeB/.style={coulCourbeB,line width=1.1pt,smooth},
  courbeC/.style={coulCourbeC,line width=1.1pt,smooth},
}

\newlist{etapes}{enumerate}{1}
\setlist[etapes]{label=\textbf{\arabic*.},leftmargin=*,itemsep=1pt,topsep=2pt}

\newlist{expressions}{enumerate}{1}
\setlist[expressions]{label={\Alph*\ $=$},leftmargin=*,itemsep=3pt,topsep=3pt}

\newcounter{exo}
\renewcommand{\theexo}{\ifnum\value{exo}<10 0\fi\arabic{exo}}
\newtcolorbox{exercice}[1][]{%
  enhanced,breakable,
  colback=white,colframe=coulExo,
  boxrule=0.9pt,arc=2pt,
  left=8pt,right=8pt,top=8pt,bottom=6pt,
  coltitle=white,fonttitle=\bfseries,
  attach boxed title to top left={xshift=8pt,yshift=-\tcboxedtitleheight/2},
  boxed title style={colback=coulExo,arc=2pt,boxrule=0pt},
  phantom={\signet{2}{Exercice \arabic{exo}}},
  title={Exercice \theexo\ifx\relax#1\relax\else{} --- #1\fi}}
\newenvironment{exo}[1][\relax]{\refstepcounter{exo}\begin{exercice}[#1]}{\end{exercice}}

  \newenvironment{correction}{\par}{\par}
  \newcommand{\autableau}{}
  \newcommand{\etapeexemples}{\etape{Exemples corrigés}}
  \newcommand{\etapetableau}[1]{\etape{#1}}
  \newenvironment{exempletableau}[1][Exemple]%
    {\begin{exemplebox}[#1]}{\end{exemplebox}}

\RequirePackage{comment}
\excludecomment{correctionavj}

\newcommand{\coursEntete}{}
\newcommand{\coursNiveau}{}
\newcommand{\coursTitre}{}
\newcommand{\coursSurTitre}{}
\newcommand{\coursSousTitre}{}

\newcommand{\enteteCours}[1]{\renewcommand{\coursEntete}{#1}}
\newcommand{\chapitrecours}[2]{\enteteCours{Chapitre #1 --- #2}}
\newcommand{\niveaucours}[1]{\renewcommand{\coursNiveau}{#1}}
\newcommand{\titrecours}[3][]{%
  \renewcommand{\coursTitre}{#2}%
  \renewcommand{\coursSurTitre}{#1}%
  \renewcommand{\coursSousTitre}{#3}}

\pagestyle{fancy}
\fancyhf{}
\fancyhead[L]{\footnotesize\color{coulPart}\textbf{\coursEntete}}
\fancyhead[R]{\footnotesize\color{coulPart}\textbf{\coursNiveau}}
\fancyfoot[C]{\footnotesize\thepage}
\renewcommand{\headrulewidth}{0.4pt}
\renewcommand{\headrule}{\hbox to\headwidth{\color{coulPart!40}\leaders\hrule height \headrulewidth\hfill}}

\setlength{\parskip}{2pt}

\AtBeginDocument{%
  \begin{center}
    \begin{tcolorbox}[enhanced,width=0.72\linewidth,colback=white,colframe=coulPart,
        boxrule=1.6pt,arc=3pt,halign=center,top=10pt,bottom=10pt]
      \ifx\coursSurTitre\empty
        {\Huge\bfseries\color{coulPart} \coursTitre}\\[4pt]
      \else
        {\Huge\bfseries\color{coulPart} \coursTitre}\\[3pt]
        {\Large\bfseries\color{coulPart} \coursSurTitre}\\[5pt]
      \fi
      {\normalsize \coursSousTitre}%
      
    \end{tcolorbox}
  \end{center}%
}
\makeatother

\chapitrecours{1}{Suites numériques}
\niveaucours{1\textsuperscript{re} mathématiques spécifiques}
\titrecours{SUITES NUMÉRIQUES}{Cours --- Première, mathématiques spécifiques}

\begin{document}

\begin{remarquebox}[Comment utiliser ce cours]
Chaque partie commence par ce qu'il faut \textbf{savoir} (définitions, propriétés), puis
vient ce qu'il faut \textbf{savoir faire} : une méthode, avec des
\textbf{exemples corrigés} faits en classe, puis des exercices
\og\textbf{À vous de jouer}\fg{} à chercher seul.

\end{remarquebox}

\partiecours{Notion de suite}

\begin{exemplebox}[Pour commencer]
Écrivons les multiples de $4$ dans l'ordre : $0$ ; $4$ ; $8$ ; $12$ ; $16$ ; \ldots{}
On leur donne un numéro, en commençant \textbf{à $0$} :
\[
  u_{0}=0 \qquad u_{1}=4 \qquad u_{2}=8 \qquad u_{3}=12 \qquad \ldots
\]
Chaque entier $n$ reçoit un nombre $u_{n}$ : c'est une \textbf{suite}. Ici, $u_{n}=4n$.
\end{exemplebox}

\begin{definitionbox}[Définition 1 --- Suite numérique]
Une \textbf{suite} associe un nombre $u_{n}$ à chaque entier $n=0$ ; $1$ ; $2$ ;
\ldots{} (on lit « $u$ indice $n$ »).
\begin{itemize}
  \item $u_{n}$ est le \textbf{terme de rang $n$} ; la suite se note $(u_{n})$.
  \item Le premier terme est $u_{0}$, le deuxième $u_{1}$, le troisième $u_{2}$\ldots
  \item Le terme \textbf{suivant} $u_{n}$ est $u_{n+1}$ ; le terme
        \textbf{précédent} est $u_{n-1}$.
\end{itemize}
\end{definitionbox}

\begin{definitionbox}[Définition 2 --- Suite définie de façon explicite]
Une suite est définie \textbf{de façon explicite} quand une formule donne $u_{n}$
directement à partir de $n$, par exemple $u_{n}=4n$. On calcule alors \textbf{n'importe
quel terme} sans connaître les précédents.
\end{definitionbox}

\begin{methodebox}[Méthode 1 --- Calculer un terme d'une suite définie de façon explicite]
\etapeexemples
\begin{questions}[label=\alph*)]
  \item Soit $u_{n}=3n+4$. Calculer $u_{0}$, $u_{1}$, $u_{2}$ et $u_{15}$.
  \item Soit $v_{n}=4n^{2}-1$. Calculer les quatre premiers termes.
  \item Une carte de bus coûte $8$~€, puis chaque trajet $1{,}50$~€. Pour $n$ trajets,
        on paie $c_{n}=8+1{,}5n$ euros. Calculer $c_{20}$ et interpréter.
\end{questions}
\begin{correction}
\cas{a)}
\begin{calculs}
u_{0} &= 3\times 0+4 &\qquad& \rem{on remplace $n$ par $0$}\\
u_{0} &= 4 &&\\
u_{1} &= 3\times 1+4=7 &&\\
u_{2} &= 3\times 2+4=10 &&\\
u_{15} &= 3\times 15+4=49 &\qquad& \rem{directement, sans les termes d'avant}
\end{calculs}
\cas{b)}
\begin{calculs}
v_{0} &= 4\times 0^{2}-1=-1 &\qquad& \rem{les quatre premiers : rangs $0$ à $3$}\\
v_{1} &= 4\times 1^{2}-1=3 &&\\
v_{2} &= 4\times 2^{2}-1=15 &\qquad& \rem{le carré d'abord : $4\times 4-1$}\\
v_{3} &= 4\times 3^{2}-1=35 &\qquad& \rem{$4\times 9-1$}
\end{calculs}
\cas{c)}
\begin{calculs}
c_{20} &= 8+1{,}5\times 20 &\qquad& \rem{on remplace $n$ par $20$}\\
c_{20} &= 8+30 &&\\
c_{20} &= 38 &&
\end{calculs}
\explique{Avec $20$ trajets, on paie $38$~€ en tout, carte comprise.}{on répond dans le
contexte}
\end{correction}

\etape{À vous de jouer}
\begin{questions}
  \item Soit $u_{n}=4n-7$. Calculer $u_{0}$, $u_{1}$ et $u_{30}$.
  \item Soit $v_{n}=n^{2}+3n$. Calculer les quatre premiers termes.
  \item Une salle de sport demande $30$~€ d'inscription, puis $6$~€ par séance. Pour
        $n$ séances, on paie $c_{n}=30+6n$ euros. Calculer $c_{12}$ et interpréter.
\end{questions}
\reponse{$u_{0}=-7$ ; $u_{1}=-3$ ; $u_{30}=113$ ;\; $v$ : $0$ ; $4$ ; $10$ ; $18$ ;\;
$c_{12}=102$ : $12$ séances coûtent $102$~€.}
\end{methodebox}

\begin{remarquebox}[Ne pas confondre $u_{n+1}$ et $u_{n}+1$]
$u_{n+1}$ est le terme \textbf{suivant} ; $u_{n}+1$ est le terme $u_{n}$ augmenté de $1$.
Avec $u_{n}=3n+4$ et $n=2$ : $u_{3}=13$, alors que $u_{2}+1=11$.
\end{remarquebox}

\afaire{exercices 1 à 5}

\partiecours{Suites définies par récurrence}

\begin{definitionbox}[Définition 3 --- Suite définie par récurrence]
Une suite est définie \textbf{par récurrence} quand on connaît :
\begin{itemize}
  \item son \textbf{premier terme} (en général $u_{0}$) ;
  \item une \textbf{relation} qui donne chaque terme à partir du \textbf{précédent} :
        $u_{n+1}$ en fonction de $u_{n}$.
\end{itemize}
Exemple : $u_{0}=6$ et $u_{n+1}=2u_{n}$ se lit « on part de $6$, et chaque terme est
le double du précédent ».
\end{definitionbox}

\begin{methodebox}[Méthode 2 --- Calculer les termes d'une suite définie par récurrence]
\etapeexemples
\begin{questions}[label=\alph*)]
  \item Soit $u_{0}=2$ et $u_{n+1}=5u_{n}$. Calculer $u_{1}$, $u_{2}$ et $u_{3}$.
  \item Soit $v_{0}=4$ et $v_{n+1}=2v_{n}-3$. Calculer $v_{1}$, $v_{2}$ et $v_{3}$.
  \item Léa a $45$ abonnés sur son compte de photos et en gagne $8$ chaque semaine. On
        note $a_{n}$ son nombre d'abonnés après $n$ semaines. Donner $a_{0}$, exprimer
        $a_{n+1}$ en fonction de $a_{n}$, puis calculer $a_{3}$.
\end{questions}
\begin{correction}
\cas{a)}
\begin{calculs}
u_{1} &= 5\times u_{0} &\qquad& \rem{la relation avec $n=0$}\\
u_{1} &= 5\times 2=10 &\qquad& \rem{on remplace $u_{0}$ par sa valeur}\\
u_{2} &= 5\times 10=50 &\qquad& \rem{$u_{2}=5\times u_{1}$}\\
u_{3} &= 5\times 50=250 &\qquad& \rem{$u_{3}=5\times u_{2}$}
\end{calculs}
\cas{b)}
\begin{calculs}
v_{1} &= 2\times 4-3=5 &\qquad& \rem{$v_{1}=2v_{0}-3$}\\
v_{2} &= 2\times 5-3=7 &\qquad& \rem{toujours avec le terme d'avant}\\
v_{3} &= 2\times 7-3=11 &&
\end{calculs}
\cas{c)}
\explique{$a_{0}=45$}{le nombre de départ}
\explique{$a_{n+1}=a_{n}+8$}{d'une semaine à la suivante, on ajoute $8$}
\explique{$a_{1}=53$ ; \ $a_{2}=61$ ; \ $a_{3}=69$.}{après $3$ semaines, Léa a $69$
abonnés}
\end{correction}

\etape{À vous de jouer}
\begin{questions}
  \item Soit $u_{0}=1$ et $u_{n+1}=3u_{n}+2$. Calculer $u_{1}$, $u_{2}$ et $u_{3}$.
  \item Soit $v_{0}=80$ et $v_{n+1}=v_{n}-12$. Calculer $v_{1}$, $v_{2}$ et $v_{3}$.
  \item À 8~h, un parking a $300$ places libres ; chaque quart d'heure, $25$ voitures
        s'y garent. On note $p_{n}$ le nombre de places libres après $n$ quarts
        d'heure. Donner $p_{0}$, exprimer $p_{n+1}$ en fonction de $p_{n}$, puis
        calculer $p_{3}$.
\end{questions}
\reponse{$u_{1}=5$ ; $u_{2}=17$ ; $u_{3}=53$ ;\; $v_{1}=68$ ; $v_{2}=56$ ; $v_{3}=44$ ;\;
$p_{0}=300$, $p_{n+1}=p_{n}-25$ et $p_{3}=225$.}
\end{methodebox}

\begin{remarquebox}[Au tableur]
Pour une suite définie par récurrence, $u_{12}$ demande d'abord $u_{11}$, donc
$u_{10}$\ldots{} jusqu'à $u_{0}$ : douze calculs, qu'un tableur fait à notre place.
On écrit le premier terme en B2, puis en B3 une formule qui utilise la cellule du dessus,
recopiée vers le bas. Pour $v_{n+1}=2v_{n}-3$ : \texttt{=2*B2-3}. L'épreuve demande
parfois cette formule.
\end{remarquebox}

\afaire{exercices 6 à 11}

\partiecours{Représentation graphique d'une suite}

\begin{definitionbox}[Définition 4 --- Représentation graphique]
Dans un repère, une suite $(u_{n})$ se représente par les \textbf{points} de
coordonnées $(n\,;u_{n})$.

On ne relie \textbf{pas} les points : $n$ ne prend que des valeurs entières.
\end{definitionbox}

\begin{methodebox}[Méthode 3 --- Représenter graphiquement une suite]
\etapeexemples
Représenter les six premiers termes de la suite définie par $u_{n}=\dfrac{n^{2}}{2}-2$.
\begin{correction}
\cas{Résolution}\par\nopagebreak
\begin{minipage}[c]{0.58\linewidth}\raggedright
\explique{On calcule les termes de rang $0$ à $5$ :}{tableau de valeurs}
\smallskip
\renewcommand{\arraystretch}{1.2}
\begin{tabular}{|c|c|c|c|c|c|c|}
\hline
$n$ & $0$ & $1$ & $2$ & $3$ & $4$ & $5$\\
\hline
$u_{n}$ & $-2$ & $-1{,}5$ & $0$ & $2{,}5$ & $6$ & $10{,}5$\\
\hline
\end{tabular}

\smallskip
\explique{On place les points $(n\,;u_{n})$, sans les relier.}{$n$ en abscisse,
$u_{n}$ en ordonnée}
\end{minipage}\hfill
\begin{minipage}[c]{0.38\linewidth}\centering
\begin{tikzpicture}
\begin{axis}[repere,
    x=0.4cm,y=0.4cm,
    xmin=-0.5,xmax=6.5,ymin=-2.5,ymax=11.5,
    xtick={1,...,6},ytick={-2,2,4,...,10},
    xlabel={$n$},ylabel={$u_{n}$}]
  \addplot[only marks,mark=x,mark options={line width=0.6pt},mark size=2pt,coulCourbe]
    coordinates {(0,-2)(1,-1.5)(2,0)(3,2.5)(4,6)(5,10.5)};
\end{axis}
\end{tikzpicture}
\end{minipage}
\end{correction}

\etape{À vous de jouer}
Représenter les cinq premiers termes de chaque suite :
\begin{questions}
  \item $v_{n}=7-2n$ ;
  \item $w_{n}=n^{2}-4$.
\end{questions}
\reponse{$v$ : $7$ ; $5$ ; $3$ ; $1$ ; $-1$ (points alignés) ;\; $w$ : $-4$ ; $-3$ ;
$0$ ; $5$ ; $12$.}
\end{methodebox}

\afaire{exercices 12 à 15}

\partiecours{Suites arithmétiques : définition}

\begin{exemplebox}[Pour commencer]
On part de $2$ et on ajoute $6$ à chaque étape : $2$ ; $8$ ; $14$ ; $20$ ; \ldots{}
C'est la suite définie par $u_{0}=2$ et $u_{n+1}=u_{n}+6$.

On part de $10$ et on retire $3$ : $10$ ; $7$ ; $4$ ; $1$ ; \ldots{} C'est la suite
définie par $v_{0}=10$ et $v_{n+1}=v_{n}-3$.
\end{exemplebox}

\begin{definitionbox}[Définition 5 --- Suite arithmétique]
Une suite $(u_{n})$ est \textbf{arithmétique} quand on passe d'un terme au suivant en
\textbf{ajoutant toujours le même nombre} $r$ : pour tout $n$,
\[
  \boxed{u_{n+1}=u_{n}+r}
\]
Le nombre $r$ s'appelle la \textbf{raison} de la suite.
\end{definitionbox}

\begin{proprietebox}[Propriété 1 --- Reconnaître une suite arithmétique]
Une suite est arithmétique quand la \textbf{différence} entre deux termes consécutifs
est toujours la même : $u_{n+1}-u_{n}=r$. Dans un énoncé : une grandeur augmente (ou
diminue) \textbf{de la même quantité} à chaque étape.
\end{proprietebox}

\begin{methodebox}[Méthode 4 --- Reconnaître et modéliser une suite arithmétique]
\etapeexemples
\begin{questions}[label=\alph*)]
  \item En 2020, une commune compte $800$ arbres en ville et en plante $150$ de plus
        chaque année. On note $u_{n}$ le nombre d'arbres en $2020+n$. Calculer $u_{1}$
        et $u_{2}$, exprimer $u_{n+1}$ en fonction de $u_{n}$, puis donner la nature de
        la suite.
  \item Le 1\textsuperscript{er} mars, une station de ski a $120$~cm de neige ; la
        hauteur baisse de $4$~cm par jour. On note $h_{n}$ la hauteur, en cm, $n$ jours
        après le 1\textsuperscript{er} mars. Modéliser par une suite.
  \item Les listes $7$ ; $11$ ; $15$ ; $19$ \;et\; $2$ ; $5$ ; $9$ ; $14$ sont-elles
        les premiers termes de suites arithmétiques ?
\end{questions}
\begin{correction}
\cas{a)}
\explique{$u_{0}=800$ ; \ $u_{1}=800+150=950$ ; \ $u_{2}=950+150=1\,100$.}{en 2022,
la commune compte $1\,100$ arbres}
\explique{$u_{n+1}=u_{n}+150$}{chaque année, on ajoute $150$}
\explique{$(u_{n})$ est arithmétique, de premier terme $u_{0}=800$ et de raison
$r=150$.}{même nombre ajouté à chaque étape}
\cas{b)}
\explique{$h_{0}=120$ et $h_{n+1}=h_{n}-4$.}{on retire $4$ chaque jour}
\explique{$(h_{n})$ est arithmétique, de raison $r=-4$.}{retirer $4$, c'est ajouter
$-4$}
\cas{c)}
\explique{$11-7=4$ ; \ $15-11=4$ ; \ $19-15=4$.}{différences successives}
\explique{Première liste : arithmétique, de raison $4$.}{les différences sont égales}
\explique{$5-2=3$ mais $9-5=4$.}{}
\explique{Deuxième liste : pas arithmétique.}{les différences changent}
\end{correction}

\etape{À vous de jouer}
\begin{questions}
  \item En 2024, une association compte $65$ bénévoles et en accueille $9$ de plus
        chaque année. On note $b_{n}$ le nombre de bénévoles en $2024+n$. Exprimer
        $b_{n+1}$ en fonction de $b_{n}$ et donner la nature de la suite.
  \item Les listes $30$ ; $26$ ; $22$ ; $18$ \;et\; $2$ ; $4$ ; $8$ ; $16$ sont-elles
        les premiers termes de suites arithmétiques ?
\end{questions}
\reponse{$b_{n+1}=b_{n}+9$ : arithmétique de raison $9$ ;\; la première oui (raison
$-4$), la seconde non ($2$, puis $4$).}
\end{methodebox}

\afaire{exercices 16 à 19}

\partiecours{Suites arithmétiques : expression en fonction de $n$}

\begin{proprietebox}[Propriété 2 --- Terme de rang $n$]
Si $(u_{n})$ est arithmétique de raison $r$, alors, pour tout entier $n$ :
\[
  \boxed{u_{n}=u_{0}+n\times r}
\]
En effet, pour aller de $u_{0}$ à $u_{n}$, on ajoute $n$ fois la raison :
$u_{1}=u_{0}+r$ ; $u_{2}=u_{0}+2r$ ; $u_{3}=u_{0}+3r$\ldots
\end{proprietebox}

\begin{methodebox}[Méthode 5 --- Exprimer une suite arithmétique en fonction de $n$]
\etapeexemples
\begin{questions}[label=\alph*)]
  \item Soit $u_{0}=12$ et $u_{n+1}=u_{n}-5$. Exprimer $u_{n}$ en fonction de $n$, puis
        calculer $u_{30}$.
  \item Soit $u_{1}=8$ et $u_{n+1}=u_{n}+3$. Exprimer $u_{n}$ en fonction de $n$.
  \item Arbres de la méthode 4 : exprimer $u_{n}$ en fonction de $n$, puis calculer le
        nombre d'arbres en 2035.
\end{questions}
\begin{correction}
\cas{a)}
\begin{calculs}
u_{n} &= u_{0}+n\times r &\qquad& \rem{premier terme $12$, raison $-5$}\\
u_{n} &= 12+n\times(-5) &&\\
u_{n} &= 12-5n &&\\
u_{30} &= 12-5\times 30 &\qquad& \rem{on remplace $n$ par $30$}\\
u_{30} &= -138 &&
\end{calculs}
\cas{b)}
\begin{calculs}
u_{0} &= u_{1}-3 &\qquad& \rem{on recule d'un rang : on retire la raison}\\
u_{0} &= 5 &&\\
u_{n} &= 5+3n &\qquad& \rem{$u_{0}+n\times r$}
\end{calculs}
\explique{Contrôle : $u_{1}=5+3\times 1=8$.}{on retrouve la donnée}
\cas{c)}
\begin{calculs}
u_{n} &= 800+150n &\qquad& \rem{premier terme $800$, raison $150$}\\
u_{15} &= 800+150\times 15 &\qquad& \rem{2035 correspond à $n=15$}\\
u_{15} &= 3\,050 &&
\end{calculs}
\explique{En 2035, la commune comptera $3\,050$ arbres.}{}
\end{correction}

\etape{À vous de jouer}
\begin{questions}
  \item Soit $v_{0}=-3$ et $v_{n+1}=v_{n}+7$. Exprimer $v_{n}$ en fonction de $n$ et
        calculer $v_{20}$.
  \item Soit $w_{1}=40$ et $w_{n+1}=w_{n}-1{,}5$. Exprimer $w_{n}$ en fonction de $n$.
  \item Bénévoles de la méthode 4 : combien seront-ils en 2034 ?
\end{questions}
\reponse{$v_{n}=-3+7n$ et $v_{20}=137$ ;\; $w_{0}=41{,}5$ donc $w_{n}=41{,}5-1{,}5n$ ;\;
$b_{n}=65+9n$ donc $b_{10}=155$.}
\end{methodebox}

\afaire{exercices 20 à 24}

\partiecours{Suites arithmétiques : problème de seuil}

\begin{remarquebox}[Problème de seuil]
« À partir de quel rang $u_{n}$ dépasse-t-il telle valeur ? » Pour une suite
arithmétique, on résout une \textbf{inéquation du premier degré} (vue en seconde).
Comme $n$ est entier, on garde le \textbf{premier entier} qui convient.
\end{remarquebox}

\begin{methodebox}[Méthode 6 --- Résoudre un problème de seuil (croissance linéaire)]
\etapeexemples
Une cycliste parcourt $20$~km le jour $0$, puis $4$~km de plus chaque jour. On note
$u_{n}$ la distance parcourue le jour $n$. À partir de quel jour parcourra-t-elle plus
de $75$~km ?
\begin{correction}
\cas{Résolution}\par\nopagebreak
\explique{$u_{n}=20+4n$}{arithmétique : premier terme $20$, raison $4$}
\begin{calculs}
20+4n &> 75 &\qquad& \rem{on traduit la question}\\
4n &> 55 &\qquad& \rem{on retire $20$}\\
n &> 13{,}75 &\qquad& \rem{on divise par $4$, positif}
\end{calculs}
\explique{Le premier entier qui convient est $n=14$ : à partir du
$14^{\text{e}}$ jour, elle parcourra plus de $75$~km.}{on répond dans le contexte}
\explique{Contrôle : $u_{13}=72$ et $u_{14}=76$.}{}
\end{correction}

\etape{À vous de jouer}
\begin{questions}
  \item Arbres de la méthode 4 ($u_{n}=800+150n$) : à partir de quelle année la commune
        comptera-t-elle au moins $2\,750$ arbres ?
  \item Neige de la méthode 4 ($h_{n}=120-4n$) : à partir de quel jour de mars y
        aura-t-il moins de $50$~cm de neige ?
\end{questions}
\reponse{$150n\ge 1\,950$ donc $n\ge 13$ : en 2033 ;\; $-4n<-70$ donc $n>17{,}5$ (on
divise par $-4$ : le sens change) : $n=18$, le 19 mars.}
\end{methodebox}

\afaire{exercices 25 à 28}

\partiecours{Suites arithmétiques : variations et représentation graphique}

\begin{proprietebox}[Propriété 3 --- Sens de variation]
Soit $(u_{n})$ une suite arithmétique de raison $r$.
\begin{itemize}
  \item Si $r>0$, la suite est \textbf{croissante} : chaque terme est plus grand que le
        précédent (les points montent).
  \item Si $r=0$, la suite est \textbf{constante}.
  \item Si $r<0$, la suite est \textbf{décroissante} : chaque terme est plus petit que
        le précédent (les points descendent).
\end{itemize}
\end{proprietebox}

\demo{$u_{n+1}=u_{n}+r$ : on ajoute $r$ à chaque terme ; si $r>0$, le terme suivant est
plus grand.}

\begin{proprietebox}[Propriété 4 --- Représentation graphique]
Les points d'une suite arithmétique sont \textbf{alignés}, sur la droite d'équation
$y=u_{0}+rx$ : on parle de \textbf{croissance linéaire}. La raison $r$ joue le rôle du
coefficient directeur.
\end{proprietebox}

\begin{center}
\begin{minipage}{0.5\linewidth}
\centering
\begin{tikzpicture}
\begin{axis}[repere,
    x=0.6cm,y=0.6cm,
    xmin=-0.5,xmax=9.5,ymin=-0.5,ymax=5.5,
    xtick={1,...,9},ytick={1,...,5},
    xlabel={$n$},ylabel={$u_{n}$}]
  \addplot[dashed,black!40,domain=0:8.6]{1+0.5*x};
  \addplot[only marks,mark=x,mark options={line width=0.6pt},mark size=2pt,coulCourbe]
    coordinates {(0,1)(1,1.5)(2,2)(3,2.5)(4,3)(5,3.5)(6,4)(7,4.5)(8,5)};
\end{axis}
\end{tikzpicture}
\end{minipage}\hfill
\begin{minipage}{0.46\linewidth}
Suite arithmétique de premier terme $u_{0}=1$ et de raison $r=0{,}5$ :
$u_{n}=1+0{,}5n$.

\smallskip
Les points sont alignés : à chaque pas de $1$ vers la droite, on monte de $0{,}5$.
Comme $r>0$, la suite est croissante.
\end{minipage}
\end{center}

\begin{methodebox}[Méthode 7 --- Sens de variation d'une suite arithmétique]
\etapeexemples
Donner le sens de variation de chaque suite.
\begin{questions}[label=\alph*)]
  \item $u_{n}=7-2n$
  \item $v_{0}=10$ et $v_{n+1}=v_{n}+0{,}3$
  \item La suite représentée ci-dessus : lire aussi son premier terme et sa raison.
\end{questions}
\begin{correction}
\cas{a)}
\explique{$(u_{n})$ est arithmétique de raison $-2$ : elle est décroissante.}{la
raison est le coefficient de $n$, et $-2<0$}
\cas{b)}
\explique{$(v_{n})$ est arithmétique de raison $0{,}3$ : elle est croissante.}{on
ajoute $0{,}3$ à chaque étape, et $0{,}3>0$}
\cas{c)}
\explique{$u_{0}=1$}{le premier point est $(0\,;1)$}
\explique{$r=0{,}5$ : la suite est croissante.}{d'un point au suivant, on monte de
$0{,}5$}
\end{correction}

\etape{À vous de jouer}
Donner le sens de variation de chaque suite arithmétique :
\[
  u_{n}=-5+0{,}5n \qquad\qquad v_{0}=3 \text{ et } v_{n+1}=v_{n}-1{,}5 \qquad\qquad
  w_{n}=4n-9
\]
\reponse{$r=0{,}5$ : croissante ;\; $r=-1{,}5$ : décroissante ;\; $r=4$ : croissante.}
\end{methodebox}

\afaire{exercices 29 à 31}

\partiecours{Suites géométriques : définition}

\begin{remarquebox}[Rappel de seconde --- Coefficient multiplicateur]
\begin{itemize}
  \item Augmenter de $t\,\%$, c'est multiplier par $1+\dfrac{t}{100}$ : $+3\,\%$ donne
        $\times 1{,}03$.
  \item Diminuer de $t\,\%$, c'est multiplier par $1-\dfrac{t}{100}$ : $-15\,\%$ donne
        $\times 0{,}85$.
\end{itemize}
\end{remarquebox}

\begin{exemplebox}[Pour commencer]
On place $1\,000$~€ sur un livret à $3\,\%$ par an, intérêts compris dans le capital :
chaque année, le capital est multiplié par $1{,}03$.
\[
  u_{0}=1\,000 \qquad u_{1}=1\,030 \qquad u_{2}=1\,060{,}90 \qquad \ldots
\]
C'est la suite définie par $u_{0}=1\,000$ et $u_{n+1}=1{,}03\,u_{n}$.
\end{exemplebox}

\begin{definitionbox}[Définition 6 --- Suite géométrique]
Une suite $(u_{n})$ est \textbf{géométrique} quand on passe d'un terme au suivant en
\textbf{multipliant toujours par le même nombre} $q>0$ : pour tout $n$,
\[
  \boxed{u_{n+1}=q\times u_{n}}
\]
Le nombre $q$ s'appelle la \textbf{raison} de la suite.
\end{definitionbox}

\begin{proprietebox}[Propriété 5 --- Reconnaître une suite géométrique]
Une suite à termes strictement positifs est géométrique quand le \textbf{quotient} de
deux termes consécutifs est toujours le même : $\dfrac{u_{n+1}}{u_{n}}=q$. Dans un
énoncé : une grandeur évolue \textbf{du même pourcentage} à chaque étape.
\end{proprietebox}

\begin{methodebox}[Méthode 8 --- Reconnaître et modéliser une suite géométrique]
\etapeexemples
\begin{questions}[label=\alph*)]
  \item En 2022, un lac contient $400$~kg de déchets plastiques ; grâce au nettoyage,
        cette masse baisse de $20\,\%$ par an. On note $v_{n}$ la masse, en kg, en
        $2022+n$. Calculer $v_{1}$, exprimer $v_{n+1}$ en fonction de $v_{n}$, puis
        donner la nature de la suite.
  \item Les listes $5$ ; $15$ ; $45$ ; $135$ \;et\; $5$ ; $10$ ; $15$ ; $20$ sont-elles
        les premiers termes de suites géométriques ?
\end{questions}
\begin{correction}
\cas{a)}
\explique{Baisser de $20\,\%$, c'est multiplier par $0{,}8$.}{$1-\dfrac{20}{100}=0{,}8$}
\explique{$v_{0}=400$ et $v_{1}=0{,}8\times 400=320$.}{en 2023, il reste $320$~kg}
\explique{$v_{n+1}=0{,}8\,v_{n}$ : $(v_{n})$ est géométrique, de premier terme
$v_{0}=400$ et de raison $q=0{,}8$.}{même coefficient à chaque étape}
\cas{b)}
\explique{$\dfrac{15}{5}=3$ ; \ $\dfrac{45}{15}=3$ ; \ $\dfrac{135}{45}=3$.}{quotients
successifs}
\explique{Première liste : géométrique, de raison $3$.}{les quotients sont égaux}
\explique{$\dfrac{10}{5}=2$ mais $\dfrac{15}{10}=1{,}5$.}{}
\explique{Deuxième liste : pas géométrique (elle est arithmétique, de raison $5$).}{}
\end{correction}

\etape{À vous de jouer}
\begin{questions}
  \item En 2025, une ville compte $30\,000$ habitants ; sa population augmente de
        $2\,\%$ par an. On note $p_{n}$ la population en $2025+n$. Donner $p_{0}$ et
        $p_{1}$, puis exprimer $p_{n+1}$ en fonction de $p_{n}$.
  \item Un téléphone acheté $800$~€ perd $25\,\%$ de sa valeur chaque année. Quelle est
        la raison de la suite de ses valeurs ?
  \item La liste $81$ ; $27$ ; $9$ ; $3$ est-elle géométrique ?
\end{questions}
\reponse{$p_{0}=30\,000$ ; $p_{1}=30\,600$ ; $p_{n+1}=1{,}02\,p_{n}$ ;\; $q=0{,}75$ ;\;
oui, de raison $\dfrac{1}{3}$.}
\end{methodebox}

\afaire{exercices 32 à 36}

\partiecours{Suites géométriques : expression en fonction de $n$}

\begin{proprietebox}[Propriété 6 --- Terme de rang $n$]
Si $(u_{n})$ est géométrique de raison $q$, alors, pour tout entier $n$ :
\[
  \boxed{u_{n}=u_{0}\times q^{n}}
\]
En effet, pour aller de $u_{0}$ à $u_{n}$, on multiplie $n$ fois par la raison :
$u_{1}=u_{0}\times q$ ; $u_{2}=u_{0}\times q^{2}$ ; $u_{3}=u_{0}\times q^{3}$\ldots
\end{proprietebox}

\begin{methodebox}[Méthode 9 --- Exprimer une suite géométrique en fonction de $n$]
\etapeexemples
\begin{questions}[label=\alph*)]
  \item Soit $u_{0}=5$ et $u_{n+1}=3u_{n}$. Exprimer $u_{n}$ en fonction de $n$, puis
        calculer $u_{4}$.
  \item Soit $u_{1}=12$ et $u_{n+1}=4u_{n}$. Exprimer $u_{n}$ en fonction de $n$.
  \item Livret à $3\,\%$ (partie VIII) : exprimer $u_{n}$ en fonction de $n$, puis donner le
        capital au bout de $10$ ans. On donne $1{,}03^{10}\approx 1{,}34$.
\end{questions}
\begin{correction}
\cas{a)}
\begin{calculs}
u_{n} &= u_{0}\times q^{n} &\qquad& \rem{premier terme $5$, raison $3$}\\
u_{n} &= 5\times 3^{n} &&\\
u_{4} &= 5\times 3^{4} &\qquad& \rem{la puissance d'abord : $3^{4}=81$}\\
u_{4} &= 405 &&
\end{calculs}
\cas{b)}
\begin{calculs}
u_{0} &= \dfrac{u_{1}}{4} &\qquad& \rem{on recule d'un rang : on divise par la raison}\\
u_{0} &= 3 &&\\
u_{n} &= 3\times 4^{n} &\qquad& \rem{$u_{0}\times q^{n}$}
\end{calculs}
\cas{c)}
\begin{calculs}
u_{n} &= 1\,000\times 1{,}03^{n} &\qquad& \rem{premier terme $1\,000$, raison $1{,}03$}\\
u_{10} &\approx 1\,000\times 1{,}34 &\qquad& \rem{on utilise la valeur donnée}\\
u_{10} &\approx 1\,340 &&
\end{calculs}
\explique{Au bout de $10$ ans, le capital vaut environ $1\,340$~€.}{}
\end{correction}

\etape{À vous de jouer}
\begin{questions}
  \item Soit $v_{0}=500$ et $v_{n+1}=0{,}8\,v_{n}$. Exprimer $v_{n}$ en fonction de $n$
        et calculer $v_{2}$.
  \item Soit $w_{1}=10$ et $w_{n+1}=2w_{n}$. Exprimer $w_{n}$ en fonction de $n$ et
        calculer $w_{5}$.
  \item On place $3\,000$~€ à $2\,\%$ par an. Quel est le capital au bout de $10$ ans ?
        On donne $1{,}02^{10}\approx 1{,}22$.
\end{questions}
\reponse{$v_{n}=500\times 0{,}8^{n}$ et $v_{2}=320$ ;\; $w_{0}=5$ ; $w_{n}=5\times 2^{n}$
et $w_{5}=160$ ;\; environ $3\,000\times 1{,}22=3\,660$~€.}
\end{methodebox}

\afaire{exercices 37 à 41}

\partiecours{Suites géométriques : variations et représentation graphique}

\begin{proprietebox}[Propriété 7 --- Sens de variation]
Soit $(u_{n})$ une suite géométrique de raison $q>0$ et de premier terme $u_{0}>0$.
\begin{itemize}
  \item Si $q>1$, la suite est \textbf{croissante}.
  \item Si $q=1$, la suite est \textbf{constante}.
  \item Si $0<q<1$, la suite est \textbf{décroissante}.
\end{itemize}
\end{proprietebox}

\begin{proprietebox}[Propriété 8 --- Représentation graphique]
Les points d'une suite géométrique \textbf{ne sont pas alignés}. Quand $q>1$, l'écart
entre deux termes consécutifs grandit de plus en plus : on parle de \textbf{croissance
exponentielle} (de \textbf{décroissance exponentielle} quand $0<q<1$).
\end{proprietebox}

\begin{center}
\begin{minipage}[t]{0.47\linewidth}
\centering {\small $u_{n}=0{,}5\times 2^{n}$ : $q=2>1$}\par\vspace{2pt}
\begin{tikzpicture}
\begin{axis}[repere,
    x=0.35cm,y=0.35cm,
    xmin=-0.5,xmax=6.5,ymin=-1,ymax=17,
    xtick={1,...,6},ytick={2,4,...,16},
    tick label style={font=\tiny},
    xlabel={$n$},ylabel={$u_{n}$}]
  \addplot[only marks,mark=x,mark options={line width=0.6pt},mark size=2pt,coulCourbe]
    coordinates {(0,0.5)(1,1)(2,2)(3,4)(4,8)(5,16)};
\end{axis}
\end{tikzpicture}
\end{minipage}\hfill
\begin{minipage}[t]{0.47\linewidth}
\centering {\small $v_{n}=16\times 0{,}5^{n}$ : $q=0{,}5<1$}\par\vspace{2pt}
\begin{tikzpicture}
\begin{axis}[repere,
    x=0.35cm,y=0.35cm,
    xmin=-0.5,xmax=6.5,ymin=-1,ymax=17,
    xtick={1,...,6},ytick={2,4,...,16},
    tick label style={font=\tiny},
    xlabel={$n$},ylabel={$v_{n}$}]
  \addplot[only marks,mark=x,mark options={line width=0.6pt},mark size=2pt,coulCourbe]
    coordinates {(0,16)(1,8)(2,4)(3,2)(4,1)(5,0.5)};
\end{axis}
\end{tikzpicture}
\end{minipage}
\end{center}

\begin{methodebox}[Méthode 10 --- Sens de variation d'une suite géométrique]
\etapeexemples
Donner le sens de variation de chaque suite.
\begin{questions}[label=\alph*)]
  \item $u_{n}=7\times 1{,}5^{n}$
  \item $v_{0}=40$ et $v_{n+1}=0{,}25\,v_{n}$
  \item La masse de déchets du lac (méthode 8 : $-20\,\%$ par an)
\end{questions}
\begin{correction}
\cas{a)}
\explique{$q=1{,}5>1$ : $(u_{n})$ est croissante.}{la raison est le nombre élevé à la
puissance $n$}
\cas{b)}
\explique{$q=0{,}25$ et $0<q<1$ : $(v_{n})$ est décroissante.}{on multiplie par
$0{,}25$ à chaque étape}
\cas{c)}
\explique{$q=0{,}8$ et $0<q<1$ : la masse de déchets diminue d'année en année.}{baisse
de $20\,\%$ : $\times 0{,}8$}
\end{correction}

\etape{À vous de jouer}
Donner le sens de variation de chaque suite géométrique :
\[
  u_{n}=50\times 0{,}9^{n} \qquad\qquad v_{0}=3 \text{ et } v_{n+1}=1{,}1\,v_{n}
  \qquad\qquad \text{une population qui augmente de } 2\,\% \text{ par an}
\]
\reponse{$q=0{,}9<1$ : décroissante ;\; $q=1{,}1>1$ : croissante ;\; $q=1{,}02>1$ :
croissante.}
\end{methodebox}

\afaire{exercices 42 à 44}

\partiecours{Suites géométriques : problème de seuil}

\begin{remarquebox}[Seuil et croissance exponentielle]
Cette année, l'inéquation $u_{0}\times q^{n}>\ldots$ ne se résout pas par le calcul.
L'énoncé donne un \textbf{tableau de valeurs} de $q^{n}$ (comme à l'épreuve).
\end{remarquebox}

\begin{methodebox}[Méthode 11 --- Résoudre un problème de seuil avec un tableau de valeurs]
\etapeexemples
Déchets du lac (méthode 8) : $v_{n}=400\times 0{,}8^{n}$. En quelle année restera-t-il
au plus $100$~kg de déchets ? On donne :
\begin{center}
\renewcommand{\arraystretch}{1.3}
\begin{tabular}{|c|c|c|c|c|}
\hline
$n$ & $5$ & $6$ & $7$ & $8$\\
\hline
$0{,}8^{n}$ (arrondi) & $0{,}33$ & $0{,}26$ & $0{,}21$ & $0{,}17$\\
\hline
\end{tabular}
\end{center}
\begin{correction}
\cas{Résolution}\par\nopagebreak
\begin{calculs}
400\times 0{,}8^{n} &\le 100 &\qquad& \rem{on traduit la question : $v_{n}\le 100$}\\
0{,}8^{n} &\le 0{,}25 &\qquad& \rem{on divise par $400$ : $\dfrac{100}{400}=0{,}25$}
\end{calculs}
\explique{$0{,}8^{6}\approx 0{,}26>0{,}25$ et $0{,}8^{7}\approx 0{,}21\le 0{,}25$ : le
premier rang est $n=7$.}{on lit le tableau}
\explique{Il restera au plus $100$~kg de déchets en $2022+7=2029$.}{on répond dans le
contexte}
\end{correction}

\etape{À vous de jouer}
Une voiture achetée $20\,000$~€ perd $15\,\%$ de sa valeur chaque année. On note $a_{n}$
sa valeur, en euros, après $n$ années.
\begin{questions}
  \item Exprimer $a_{n}$ en fonction de $n$.
  \item Au bout de combien d'années vaudra-t-elle moins de $10\,000$~€ ? On donne
        $0{,}85^{3}\approx 0{,}61$ ; $0{,}85^{4}\approx 0{,}52$ ;
        $0{,}85^{5}\approx 0{,}44$.
\end{questions}
\reponse{$a_{n}=20\,000\times 0{,}85^{n}$ ;\; $a_{n}<10\,000$ revient à
$0{,}85^{n}<0{,}5$ : au bout de $5$ ans.}
\end{methodebox}

\afaire{exercices 45 à 47}

\besoinplace{20\baselineskip}
\partiecours{Comparer deux évolutions}

\begin{center}
\renewcommand{\arraystretch}{1.5}
\begin{tabular}{|l|c|c|}
\hline
 & \textbf{Suite arithmétique} & \textbf{Suite géométrique}\\
\hline
D'un terme au suivant & on \textbf{ajoute} $r$ & on \textbf{multiplie} par $q$\\
\hline
Relation de récurrence & $u_{n+1}=u_{n}+r$ & $u_{n+1}=q\times u_{n}$\\
\hline
Terme de rang $n$ & $u_{n}=u_{0}+nr$ & $u_{n}=u_{0}\times q^{n}$\\
\hline
Croissante quand & $r>0$ & $q>1$\\
\hline
Décroissante quand & $r<0$ & $0<q<1$\\
\hline
Graphique & points alignés & points non alignés\\
\hline
Dans un énoncé & même quantité ajoutée & même pourcentage d'évolution\\
\hline
Croissance & linéaire & exponentielle\\
\hline
\end{tabular}
\end{center}

\begin{methodebox}[Méthode 12 --- Comparer une croissance linéaire et une croissance exponentielle]
\etapeexemples
Au 1\textsuperscript{er} janvier, deux chaînes vidéo ont chacune $5\,000$ abonnés. La
chaîne A en gagne $300$ par mois ; celle de la chaîne B augmente de $5\,\%$ par mois. On
note $u_{n}$ et $v_{n}$ les nombres d'abonnés de A et de B après $n$ mois. À partir de
quel mois la chaîne B a-t-elle plus d'abonnés que la chaîne A ? On donne :
\begin{center}
\renewcommand{\arraystretch}{1.3}
\begin{tabular}{|c|c|c|c|c|}
\hline
$n$ & $7$ & $8$ & $9$ & $10$\\
\hline
$v_{n}$ (arrondi) & $7\,036$ & $7\,387$ & $7\,757$ & $8\,144$\\
\hline
\end{tabular}
\end{center}
\begin{correction}
\cas{Résolution}\par\nopagebreak
\explique{$u_{n}=5\,000+300n$}{A : on ajoute $300$, suite arithmétique}
\explique{$v_{n}=5\,000\times 1{,}05^{n}$}{B : on multiplie par $1{,}05$, suite
géométrique}
\begin{calculs}
u_{8} &= 5\,000+300\times 8=7\,400 &\qquad& \rem{$7\,400>7\,387$ : A devant}\\
u_{9} &= 5\,000+300\times 9=7\,700 &\qquad& \rem{$7\,700<7\,757$ : B devant}
\end{calculs}
\explique{À partir du $9^{\text{e}}$ mois, la chaîne B a plus d'abonnés.}{la croissance
exponentielle finit par dépasser la croissance linéaire}
\end{correction}

\etape{À vous de jouer}
En 2025, la ville A compte $12\,000$ habitants et en gagne $400$ par an ; la ville B
compte $10\,000$ habitants, avec une hausse de $4\,\%$ par an.
\begin{questions}
  \item Exprimer la population de chaque ville en fonction de $n$ (année $2025+n$).
  \item À partir de quelle année la ville B sera-t-elle la plus peuplée ? On donne, pour
        la ville B : $16\,651$ habitants en 2038, $17\,317$ en 2039 et $18\,009$ en 2040
        (valeurs arrondies).
\end{questions}
\reponse{$a_{n}=12\,000+400n$ et $b_{n}=10\,000\times 1{,}04^{n}$ ;\; $a_{14}=17\,600>
17\,317$ mais $a_{15}=18\,000<18\,009$ : à partir de 2040.}
\end{methodebox}

\afaire{exercices 48 à 50, puis les exercices de type épreuve 51 à 53}

\end{document}
